The central empirical result is a dissociation. Chilean elite-coded surnames often produce large latent status associations, including against a rarity-matched control, but those associations usually do not propagate into matched consequential scores. This matters for evaluation design because an association probe and a decision audit support different claims. In the present data, treating the former as evidence of the latter would substantially overstate downstream bias for several models.

The dissociation is not explained by a single universal alignment strategy. Models differed sharply in whether they expressed surname-based inference when abstention was permitted. Some systems, such as Gemini 3.6 Flash and Claude Sonnet 5, selectively expressed status inferences for elite-coded surnames, while several others generally abstained. Yet both patterns could coexist with very small consequential elite-minus-common effects. The result therefore suggests that explicit willingness to infer status, latent forced association, and use of the signal in a constrained decision are at least partially separable behaviors.

Decision structure also matters, but the secondary evidence does not support a broad claim that holistic or metadata presentation systematically unlocks elite-status leakage. No elite-minus-common metadata or holistic contrast survived FDR correction. The clearest secondary effect was instead Qwen 3.7 Max's general reduction in scores whenever a surname was visible, including elite, common, and rare controls. This illustrates why blind baselines and rarity controls are important: a name-presence effect can otherwise be mistaken for a class-specific effect.

Llama 4 Maverick warrants cautious attention. Its primary elite-minus-common contrast was nominally positive and corresponded to 0.095 standard deviations, just below the magnitude used to define the $\pm0.10$ SD equivalence region. The result failed the equivalence test and its confidence interval excludes zero by a small amount. However, it is a small effect, it is not accompanied by strong model-level association-to-leakage coupling, and the predeclared robustness study was not executed. It should therefore be treated as a model-specific signal for replication rather than evidence of a general elite-surname advantage.

More broadly, association magnitude was a poor predictor of leakage. Gemini 3.6 Flash and Claude Sonnet 5 showed two of the strongest association contrasts while exhibiting decision effects essentially at zero. Conversely, some weaker-association systems displayed greater decision-score variability. This pattern is consistent with the construct-validity motivation of the study: internalized or elicitable social knowledge is not sufficient to infer consequential treatment under a given decision procedure.

The findings should not be read as proving that surname-mediated discrimination is absent from deployed language-model systems. The study evaluates synthetic tasks, fixed provider endpoints, a specific Chilean-Spanish instrument, and bounded decision scaffolds. Different prompting, tool use, long-horizon interaction, fine-tuning, or institutional context may alter whether latent associations become behaviorally relevant. The contribution here is narrower: it shows empirically that association and leakage can diverge strongly, and provides an auditable design for measuring that transition rather than assuming it.